\uuid{aEtH}
\titre{Écriture décimale}
\auteur{bodin ; hanani ; mzari}
\organisation{exo7}
\licence{CC BY-NC-SA 4.0}
\datecreate{2020-01-20}
\niveau{L1}
\module{Analyse}
\chapitre{Propriétés de R}
\sousChapitre{Les rationnels}
\difficulte{3}
\theme{Rationnels, écriture décimale}
\exo7qcmid{lille-1.reels.005}

\contenu{
\begin{question}
\qclassification{120.01}
\qtags{motcle=écriture décimale}
Quelles sont les assertions vraies ?
\begin{answers}
    \bad{L'écriture décimale de $\sqrt{3}$ est finie ou périodique.}

    \good{L'écriture décimale de $\frac{n}{n+1}$ est finie ou périodique (quelque soit $n\in\Nn$).}

    \bad{Un nombre réel qui admet une écriture décimale infinie est un nombre irrationnel.}

    \good{Un nombre réel qui admet une écriture décimale finie est un nombre rationnel.}

\end{answers}
\begin{explanations}
Les nombre rationnels sont exactement les nombres qui admettent une écriture décimale finie ou périodique.
\end{explanations}
\end{question}
}
